% =====================================================================
% APÊNDICES. Os arquivos apendice_*_texto.tex e apendice_tabela_respostas.tex
% são GERADOS a partir do repositório (prompts, banco curado e respostas da
% primeira rodada) e não devem ser editados à mão.
% =====================================================================

\chapter{Prompts utilizados no sistema}
\label{ap:prompts}

Este apêndice reproduz os dois \emph{prompts} do sistema na versão usada para gerar o banco curado.
Eles são artefatos versionados da pesquisa e ficam no repositório em
\texttt{sistema/prompts/}. O Verificador Simbólico e o Orquestrador não têm \emph{prompt}, porque
não usam modelo de linguagem. Além do texto fixo reproduzido a seguir, cada pedido ao Gerador inclui
a especificação do professor, os descritores operacionais de dificuldade e natureza e as exigências
da habilidade, e cada pedido ao Crítico inclui a questão, a especificação e as mesmas exigências.
As linhas mais longas foram quebradas para caber na página.

% Gerado a partir de sistema/prompts/*.md. Não editar à mão: regerar se os prompts mudarem.
\section*{Prompt do Agente Gerador}
\addcontentsline{toc}{section}{Prompt do Agente Gerador}
{\footnotesize
\begin{verbatim}
Você é um elaborador experiente de questões de Matemática para o Ensino Médio
  brasileiro,
que segue a BNCC e as boas práticas de elaboração de itens (clareza do enunciado,
  caminho
de solução não trivial, resposta determinada e verificável).

Sua tarefa: produzir UMA questão conforme a especificação fornecida pelo professor.

REGRAS OBRIGATÓRIAS:
1. Responda APENAS com um objeto JSON válido, sem texto antes ou depois, sem cercas
  de código.
2. O enunciado deve ser claro, sem ambiguidade, com todos os dados necessários e
  nenhum dado
   supérfluo não intencional. Não sinalize o procedimento de resolução no enunciado.
3. A resolução deve ser passo a passo, em Markdown, com notação LaTeX entre $...$.
4. Em múltipla escolha: exatamente 4 alternativas, uma correta; cada distrator deve
   corresponder a um erro sistemático plausível de estudante (declare qual erro).
   O sistema sorteia a ordem das alternativas depois de você: no gabarito e na
     resolução,
   identifique a resposta pelo conteúdo, nunca pela letra ou pela posição.
5. Preencha "verificaveis" com a formalização SymPy do problema. É uma LISTA: inclua
  uma
   entrada para CADA afirmação que a questão faz. Uma questão que articula dois
     temas tem
   duas; uma que pede um valor e uma propriedade também.
   - tipo "equacao": expressao = "Eq(...)", resposta_esperada = lista de soluções,
     ex.: "[1, Rational(3,2)]"
   - tipo "funcao": expressao = f(x); parametros.consulta em
     {"zeros","vertice","valor","maximo","minimo","dominio","imagem","periodo","cres
       cimento"}.
     Para "valor", inclua parametros.ponto. resposta_esperada conforme a consulta:
       zeros -> lista de raízes | vertice -> "[xv, yv]" | valor/maximo/minimo ->
         valor único
       dominio, imagem -> conjunto em sintaxe SymPy: "Interval.open(0, oo)",
         "Interval(-2, 4)",
         "S.Reals", "S.Naturals". Em "dominio", declare o domínio COMO A QUESTÃO O
           DEFINE
         (se t é tempo, "Interval(0, oo)"; se n conta termos, "S.Naturals"). Em
           "imagem",
         inclua parametros.dominio com esse mesmo conjunto quando a questão
           restringir o domínio.
       periodo -> o período fundamental, ex.: "2*pi" ou "pi"
       crescimento -> exatamente a palavra "crescente" ou "decrescente"
     Escreva logaritmo como log(x) (natural) ou log(x, 2) (base 2); seno e cosseno
       como
     sin(x) e cos(x). Função por várias sentenças: Piecewise((expr1, cond1), (expr2,
       True)).
   - tipo "progressao": parametros = {tipo_progressao: "pa"|"pg", a1, razao, n,
     consulta: "termo"|"soma"};
     resposta_esperada = valor único. expressao pode ser "-".
   - tipo "propriedade": use quando o que se verifica NÃO é um valor, mas uma
     expressão que o
     estudante deve produzir -- generalizar um padrão de uma tabela, ou associar uma
       progressão
     à função correspondente. resposta_esperada = a expressão (ex.: "3*n + 2");
       expressao = "-";
     parametros com um ou mais de:
       "pontos": "[(1,5),(2,8),(3,11)]"  -- a expressão deve valer y_i em cada x_i
       "grau": "1"                        -- grau do polinômio
       "forma": "a*x**2"                  -- só os termos dessa forma (nenhum de
         grau menor)
       "sequencia": "pa"|"pg" com "a1" e "razao" -- a expressão deve coincidir com a
         progressão
     Declare TODOS os predicados que a questão sustenta: cada um é uma conferência a
       mais.
   Use sintaxe SymPy exata: Rational(a,b) para frações, sqrt() para raízes, ** para
     potência.
   A resposta_esperada DEVE ser exatamente o gabarito da questão.
6. Se a questão não puder ser formalizada nesses tipos, use "verificaveis": [] --
  mas prefira
   sempre questões formalizáveis, e formalize tudo o que der.
7. A habilidade da BNCC não é um rótulo: ela determina o que a questão precisa
  exigir do
   estudante. Cumpra as exigências listadas na especificação -- uma questão sobre o
     conteúdo
   certo que não faça o que a habilidade pede NÃO serve. Exemplo: numa habilidade de
     converter
   representação algébrica em geométrica, pedir apenas as raízes não realiza a
     habilidade,
   ainda que a matemática esteja correta.
8. Quando a especificação trouxer mais de um tema, produza UMA questão que os
  articule no mesmo
   problema (a relação entre eles é parte do que se avalia), e não dois itens
     justapostos.
9. Se receber FEEDBACK de uma tentativa anterior, corrija exatamente o que foi
  apontado,
   preservando o que estava bom.

ESQUEMA DO JSON DE SAÍDA:
{
  "enunciado": "...",
  "resolucao": "...",
  "gabarito": "...",
  "alternativas": [{"texto": "...", "correta": true|false, "erro_representado":
    "..."|null}] | null,
  "verificaveis": [
    {
      "tipo": "equacao|funcao|progressao|propriedade",
      "expressao": "...",
      "incognitas": ["x"],
      "resposta_esperada": "...",
      "parametros": {}
    }
  ]
}
\end{verbatim}
}
\section*{Prompt do Agente Crítico Didático}
\addcontentsline{toc}{section}{Prompt do Agente Crítico Didático}
{\footnotesize
\begin{verbatim}
Você é um avaliador didático de questões de Matemática do Ensino Médio, com formação
  em
educação matemática. Avalie a questão recebida segundo a rubrica abaixo, derivada de
referenciais consolidados (Pólya, Teoria das Situações Didáticas, taxonomias de
  Bloom e SOLO,
BNCC).

A correção matemática já foi verificada por um sistema simbólico -- NÃO a reavalie.
Avalie os cinco critérios seguintes, cada um com nota de 1 a 5 e comentário objetivo:

1. "clareza": O enunciado permite identificar sem dúvida o que é dado, o que é
  pedido e as
   condições? Há ambiguidade lexical, estrutural ou de contexto? Os dados são
     completos?
2. "adequacao_nivel": O processo cognitivo efetivamente exigido corresponde ao nível
  de Bloom
   declarado na especificação? A estrutura de resposta esperada (na taxonomia SOLO)
     é coerente
   com esse nível -- ex.: uma questão "analisar" não pode admitir resposta correta
     puramente
   multiestrutural? Os conteúdos são compatíveis com o Ensino Médio?
3. "alinhamento_bncc": Julgue pelas EXIGÊNCIAS listadas na especificação, uma a uma
  -- não pela
   impressão geral de que "é sobre o assunto". Abordar o conteúdo certo não basta:
     se a questão
   não faz o que a habilidade pede, a nota é 2 ou menos, ainda que a matemática
     esteja perfeita
   e o enunciado claro. Quando houver mais de um tema, verifique se a questão
     realmente os
   articula num único problema em vez de justapor itens independentes.
4. "distratores": (só para múltipla escolha; senão nota 5 com comentário "não se
  aplica")
   Cada distrator corresponde a um erro sistemático plausível de estudante? Algum é
     absurdo
   ou trivialmente eliminável?
5. "originalidade": A questão evita reproduzir mecanicamente enunciados clássicos de
  livro
   didático? O contexto é significativo (quando aplicada) e o enunciado evita o
     "efeito
   Topaze" (pistas que pavimentam a solução)?

DECISÃO: "aprovado" = true somente se NENHUM critério tiver nota menor que 3.
Se reprovar, preencha "sugestoes_revisao" com instruções concretas e acionáveis para o
elaborador corrigir a questão (o que mudar, não apenas o que está errado).

Responda APENAS com JSON válido, sem texto antes ou depois:
{
  "notas": [
    {"criterio": "clareza", "nota": 1-5, "comentario": "..."},
    {"criterio": "adequacao_nivel", "nota": 1-5, "comentario": "..."},
    {"criterio": "alinhamento_bncc", "nota": 1-5, "comentario": "..."},
    {"criterio": "distratores", "nota": 1-5, "comentario": "..."},
    {"criterio": "originalidade", "nota": 1-5, "comentario": "..."}
  ],
  "aprovado": true|false,
  "sugestoes_revisao": "..." | null
}
\end{verbatim}
}


\chapter{Exemplos de questões geradas}
\label{ap:exemplos}

As questões a seguir são as seis do bloco-âncora da primeira rodada de avaliação
(Seção~\ref{sub:rodada1-decisao}), vistas pelos cinco professores, acompanhadas da Q061 na versão
corrigida depois da avaliação (Seção~\ref{sub:rodada1-erros}). O texto é o do banco curado e foi
transcrito sem edição, com gabarito e resolução. A ordem das alternativas é a do banco, que foi
sorteada depois da avaliação e por isso pode diferir da que os professores viram. O cabeçalho de cada
questão informa o formato, a dificuldade declarada, o nível cognitivo, a natureza e a garantia
obtida. As 85 questões do banco estão no repositório, em PDF e em JSON.

% Gerado a partir de banco_curado/banco_curado.tex. Não editar à mão.
\begingroup\sloppy\setlength{\parindent}{0pt}\setlength{\parskip}{0.5em}
\bigskip\noindent\textbf{Q007 · múltipla escolha · fácil · aplicar · teórica · gabarito conferido}\par\nopagebreak
Uma oficina mecânica cobra por seus serviços de revisão um valor fixo de mão de obra, mais uma taxa que depende do número de horas de trabalho. Uma revisão que durou 2 horas custou R\$\,150,00, e outra revisão, no mesmo padrão de cobrança, durou 5 horas e custou R\$\,300,00. Sabendo que o custo total varia linearmente com o número de horas trabalhadas, qual seria o custo de uma revisão que durasse 8 horas?
\begin{enumerate}[label=(\alph*)]
\item R\$\,480,00
\item R\$\,400,00
\item R\$\,600,00
\item R\$\,450,00
\end{enumerate}

\textbf{Gabarito.} R\$\,450,00

\textbf{Resolução.} Como o custo total $C$ varia linearmente com o número de horas $x$, podemos escrever $C(x) = ax + b$, em que $a$ é a taxa cobrada por hora e $b$ é o valor fixo de mão de obra.

\textbf{Passo 1 — Determinar a taxa por hora ($a$).}
Usando os dois pares de valores dados, $(2, 150)$ e $(5, 300)$:
$$a = \frac{300 - 150}{5 - 2} = \frac{150}{3} = 50$$

Ou seja, cada hora de trabalho custa R\$\,50,00.

\textbf{Passo 2 — Determinar o valor fixo ($b$).}
Substituindo o ponto $(2, 150)$ em $C(x) = 50x + b$:
$$150 = 50 \cdot 2 + b \;\Rightarrow\; 150 = 100 + b \;\Rightarrow\; b = 50$$

Logo, o modelo do custo é:
$$C(x) = 50x + 50$$

\textbf{Passo 3 — Calcular o custo para $x = 8$ horas.}
$$C(8) = 50 \cdot 8 + 50 = 400 + 50 = 450$$

Portanto, a revisão de 8 horas custaria \textbf{R\$\,450,00}.

\bigskip\noindent\textbf{Q015 · discursiva · média · analisar · aplicada · gabarito conferido}\par\nopagebreak
Uma colônia de bactérias em um experimento de laboratório tem seu número de indivíduos, a partir do início da observação, dado por $N(t) = 100 \cdot 2^{t}$, em que $t$ é o tempo em horas ($t \geq 0$) e $N(t)$ é o número de bactérias. Um pesquisador quer saber, ao contrário, quanto tempo é necessário para que a colônia atinja um determinado número $N$ de bactérias, e por isso utiliza a função inversa de $N(t)$, que ele chama de $t(N)$.

a) Determine o domínio e a imagem de $N(t)$, considerando o contexto do experimento.

b) Encontre a expressão de $t(N)$ (a função inversa de $N(t)$) e determine seu domínio e sua imagem.

c) Compare o crescimento das duas funções: as duas são crescentes, mas uma cresce 'cada vez mais rápido' e a outra cresce 'cada vez mais devagar' à medida que a variável aumenta. Identifique qual é qual e justifique observando como varia o incremento de cada função.

d) Explique por que o domínio de $N(t)$ coincide com a imagem de $t(N)$, e a imagem de $N(t)$ coincide com o domínio de $t(N)$, relacionando esse fato com a posição dos gráficos das duas funções no plano cartesiano.

\textbf{Gabarito.} a) $D_N=[0,+\infty)$, $Im_N=[100,+\infty)$. b) $t(N)=\log_2(N/100)$, com $D_t=[100,+\infty)$, $Im_t=[0,+\infty)$. c) $N(t)$ cresce cada vez mais rápido (convexa); $t(N)$ cresce cada vez mais devagar (côncava). d) Por serem funções inversas, seus gráficos são simétricos em relação à reta $y=x$, o que troca domínio e imagem entre as duas funções.

\textbf{Resolução.} \textbf{a) Domínio e imagem de $N(t) = 100\cdot 2^t$}

Como $t$ representa tempo a partir do início da observação, $t \geq 0$. Logo:
$$D_N = [0, +\infty)$$

Quando $t=0$, $N(0) = 100 \cdot 2^0 = 100$, o menor valor possível. Como $2^t$ é crescente e ilimitada quando $t \to +\infty$, $N(t)$ também cresce sem limite. Logo:
$$Im_N = [100, +\infty)$$

\textbf{b) Função inversa $t(N)$}

De $N = 100\cdot 2^t$, isolamos $t$:
$$\frac{N}{100} = 2^t \implies t = \log_2\left(\frac{N}{100}\right)$$

Como funções inversas trocam domínio e imagem entre si:
$$D_t = Im_N = [100, +\infty), \qquad Im_t = D_N = [0, +\infty)$$

(De fato, o menor $N$ possível no experimento é 100, e para $N=100$ temos $t=\log_2(1)=0$, consistente.)

\textbf{c) Comparando o crescimento}

Ambas as funções são estritamente crescentes, pois a base $2$ é maior que $1$ tanto na exponencial quanto no logaritmo.

Porém, o \textbf{modo} de crescer é diferente. Observe os incrementos de $N(t)$ a cada hora:
\begin{itemize}[nosep,leftmargin=*]
\item de $t=0$ a $t=1$: $N$ vai de $100$ a $200$ (aumenta 100)
\item de $t=1$ a $t=2$: $N$ vai de $200$ a $400$ (aumenta 200)
\item de $t=2$ a $t=3$: $N$ vai de $400$ a $800$ (aumenta 400)
\end{itemize}

Ou seja, o incremento absoluto de $N(t)$ \textbf{aumenta} a cada intervalo — a função exponencial cresce cada vez mais rápido (é convexa, sua taxa de variação aumenta).

Já para $t(N)$: cada vez que $N$ \textbf{dobra}, $t$ aumenta sempre a mesma quantidade fixa, 1 hora (pois $\log_2(2N/100) - \log_2(N/100) = \log_2 2 = 1$). Mas para produzir esse mesmo aumento de 1 unidade em $t$, é preciso um incremento absoluto de $N$ cada vez maior (de $100\to200$, depois $200\to400$, depois $400\to800$...). Isso significa que, para incrementos iguais de $N$, o aumento de $t$ fica cada vez menor — a função logarítmica cresce cada vez mais devagar (é côncava, sua taxa de variação diminui).

\textbf{d) Relação entre domínio e imagem das duas funções}

Como $t(N)$ é a função inversa de $N(t)$, seus gráficos são simétricos em relação à reta $y=x$ no plano cartesiano. Essa simetria troca os eixos: o que era eixo do domínio (horizontal) de uma passa a ser o eixo da imagem (vertical) da outra. Por isso, necessariamente:
$$D_{N} = Im_{t} = [0,+\infty) \qquad \text{e} \qquad Im_{N} = D_{t} = [100,+\infty)$$

Essa é uma propriedade geral de qualquer par de funções inversas, e em particular da relação entre exponencial e logaritmo de mesma base: uma 'desfaz' a outra, e seus papéis de domínio e imagem se invertem.

\bigskip\noindent\textbf{Q024 · múltipla escolha · difícil · criar · teórica · gabarito conferido}\par\nopagebreak
A tabela abaixo relaciona valores de $x$ e os correspondentes valores de $y$, obtidos a partir de uma mesma regra de formação (os valores de $x$ não variam em intervalos iguais):

\begin{center}
\begin{tabular}{|p{0.172\linewidth}|p{0.172\linewidth}|p{0.172\linewidth}|p{0.172\linewidth}|p{0.172\linewidth}|}
\hline
\textbf{$x$} & \textbf{$-3$} & \textbf{$1$} & \textbf{$4$} & \textbf{$8$} \\
\hline
$y$ & $14$ & $2$ & $-7$ & $-19$ \\
\hline
\end{tabular}
\end{center}

Um estudante, analisando os pares ordenados $(x,y)$ da tabela, conjectura que essa relação pode ser representada por uma função polinomial do 1º grau. Assinale a alternativa que apresenta a lei de formação $f(x)$ que generaliza corretamente o padrão observado na tabela.
\begin{enumerate}[label=(\alph*)]
\item $f(x) = -11x + 13$
\item $f(x) = -12x + 14$
\item $f(x) = -3x + 5$
\item $f(x) = 3x - 1$
\end{enumerate}

\textbf{Gabarito.} f(x) = -3x + 5

\textbf{Resolução.} \textbf{Passo 1 — Verificar se a razão entre as variações é constante (condição para função afim).}

Como os valores de $x$ não estão igualmente espaçados, é preciso calcular a razão $\dfrac{\Delta y}{\Delta x}$ entre pares consecutivos da tabela, e não apenas $\Delta y$.

Entre $x=-3$ e $x=1$: $\Delta x = 1-(-3)=4$, $\Delta y = 2-14=-12$, logo $\dfrac{\Delta y}{\Delta x}=\dfrac{-12}{4}=-3$.

Entre $x=1$ e $x=4$: $\Delta x = 4-1=3$, $\Delta y = -7-2=-9$, logo $\dfrac{\Delta y}{\Delta x}=\dfrac{-9}{3}=-3$.

Entre $x=4$ e $x=8$: $\Delta x = 8-4=4$, $\Delta y = -19-(-7)=-12$, logo $\dfrac{\Delta y}{\Delta x}=\dfrac{-12}{4}=-3$.

\textbf{Passo 2 — Concluir o tipo de função.}

A razão $\dfrac{\Delta y}{\Delta x}$ é constante e igual a $-3$ para todos os pares consecutivos, mesmo com intervalos diferentes de $x$. Isso caracteriza uma \textbf{função polinomial do 1º grau (afim)}, da forma $f(x)=mx+b$, com $m=-3$.

\textbf{Passo 3 — Determinar o coeficiente linear $b$.}

Usando o ponto $(1,2)$: $2 = -3(1)+b \Rightarrow b = 2+3 = 5$.

Logo, $f(x) = -3x+5$.

\textbf{Passo 4 — Verificar com os demais pontos.}

$f(-3) = -3(-3)+5 = 9+5 = 14$ $\checkmark$

$f(4) = -3(4)+5 = -12+5 = -7$ $\checkmark$

$f(8) = -3(8)+5 = -24+5 = -19$ $\checkmark$

Todos os pares confirmam a lei $f(x) = -3x+5$.

\bigskip\noindent\textbf{Q034 · múltipla escolha · fácil · aplicar · aplicada · gabarito conferido}\par\nopagebreak
Duas empresas de aplicativo de transporte, A e B, calculam o valor da corrida em função da distância percorrida $x$ (em km). A empresa A cobra R\$\,2,50 por quilômetro rodado, sem nenhuma taxa adicional. A empresa B cobra uma taxa fixa de embarque de R\$\,3,00 mais R\$\,2,00 por quilômetro rodado. Sejam $y_A(x)$ e $y_B(x)$, em reais, os custos totais das corridas das empresas A e B, respectivamente, em função de $x$. Se essas duas funções forem representadas por retas em um plano cartesiano (custo $y$ no eixo vertical e distância $x$ no eixo horizontal), assinale a alternativa que descreve corretamente essas retas.
\begin{enumerate}[label=(\alph*)]
\item A reta que representa a empresa B passa pela origem (0,0), pois seu coeficiente angular é 2; a reta da empresa A corta o eixo y no ponto (0,3), pois seu coeficiente angular é 2,5.
\item Ambas as retas passam pela origem (0,0), já que representam custo por quilômetro rodado; a reta de A tem coeficiente angular 2,5 e a de B tem coeficiente angular 2.
\item A reta que representa a empresa A passa pela origem (0,0) e tem coeficiente angular 2,5, sendo uma função proporcional; a reta da empresa B corta o eixo y no ponto (0,3) e tem coeficiente angular 2, não sendo proporcional.
\item A reta da empresa A corta o eixo y no ponto (0, 2,5); a reta da empresa B corta o eixo y no ponto (0,2) e tem coeficiente angular 3.
\end{enumerate}

\textbf{Gabarito.} A reta que representa a empresa A passa pela origem $(0,0)$ e tem coeficiente angular 2,5 (função proporcional); a reta da empresa B corta o eixo y no ponto $(0,3)$ e tem coeficiente angular 2 (função afim, não proporcional).

\textbf{Resolução.} \textbf{Passo 1 — Escrever as leis das funções.}

Empresa A: cobra só por km rodado, sem taxa fixa: $y_A(x) = 2{,}5x$.

Empresa B: cobra taxa fixa de R\$\,3,00 mais R\$\,2,00 por km: $y_B(x) = 2x + 3$.

\textbf{Passo 2 — Analisar a natureza de cada função.}

$y_A(x) = 2{,}5x$ é da forma $y = ax$ (sem termo constante), logo é uma função \textbf{linear/proporcional}. Seu gráfico é uma reta que passa pela origem $(0,0)$, já que $y_A(0) = 2{,}5 \cdot 0 = 0$. O coeficiente angular é $a = 2{,}5$.

$y_B(x) = 2x+3$ é da forma $y = ax+b$ com $b = 3 \neq 0$, logo é uma função \textbf{afim não proporcional}. Seu gráfico é uma reta que \textbf{não} passa pela origem, pois $y_B(0) = 3 \neq 0$. Ela corta o eixo $y$ no ponto $(0,3)$ e tem coeficiente angular $a = 2$.

\textbf{Passo 3 — Conclusão geométrica.}

A reta de A passa pela origem (comportamento proporcional, coeficiente angular 2,5); a reta de B corta o eixo $y$ em $(0,3)$ (comportamento afim, não proporcional, coeficiente angular 2). As duas retas são crescentes, pois ambos os coeficientes angulares são positivos, mas apenas a de A representa proporcionalidade direta entre custo e distância.

\bigskip\noindent\textbf{Q065 · discursiva · difícil · aplicar · aplicada · gabarito conferido}\par\nopagebreak
Um oceanógrafo registra a altura da maré (em metros) em um pequeno porto ao longo de um dia. Ele observa que a maré varia de forma periódica e simétrica em torno de um nível médio, atingindo seu valor máximo de 1,8 m exatamente à meia-noite ($t=0$, com $t$ em horas) e seu valor mínimo de 0,2 m exatamente às 6h da manhã. Esse padrão se repete a cada 12 horas.

a) Modele a altura da maré, em metros, por uma função do tipo $h(t) = A\cos(Bt) + D$, com $t$ em horas contado a partir da meia-noite. Determine explicitamente os valores de $A$, $B$ e $D$, justificando cada um a partir dos dados do fenômeno (amplitude, período e deslocamento vertical do gráfico em relação ao eixo $t$).

b) Determine todos os horários dentro das primeiras 12 horas do dia (isto é, para $0 \le t < 12$) em que a maré atinge exatamente 1,4 m de altura. Para cada horário encontrado, indique se a maré está subindo ou descendo naquele instante, justificando sua resposta apenas a partir do comportamento gráfico da função cosseno (sem usar derivadas).

\textbf{Gabarito.} h(t) = 0,8·cos($\pi$t/6) + 1,0 (A=0,8 m; B=$\pi$/6 rad/h; D=1,0 m). A maré atinge 1,4 m às 2h (descendo) e às 10h (subindo).

\textbf{Resolução.} \textbf{a) Construção do modelo}

O valor máximo da maré é $1{,}8$ m e o mínimo é $0{,}2$ m. Como $h(t)=A\cos(Bt)+D$ oscila entre $D-A$ e $D+A$:

$$D = \frac{1{,}8+0{,}2}{2} = 1{,}0 \qquad A = \frac{1{,}8-0{,}2}{2} = 0{,}8$$

O fenômeno se repete a cada 12 horas, logo o período é $T=12$. Como $T=\dfrac{2\pi}{B}$:

$$B = \frac{2\pi}{12} = \frac{\pi}{6}$$

Como a maré atinge seu \textbf{máximo} exatamente em $t=0$, e $\cos(0)=1$ é o valor máximo do cosseno, a escolha de cosseno (sem deslocamento horizontal) é coerente com o fenômeno. Assim:

$$h(t) = 0{,}8\cos\!\left(\frac{\pi t}{6}\right) + 1{,}0$$

Verificação: $h(0)=0{,}8(1)+1=1{,}8$ (máximo, ok) e $h(6)=0{,}8\cos(\pi)+1=0{,}8(-1)+1=0{,}2$ (mínimo, ok).

\textbf{b) Horários em que $h(t)=1{,}4$}

$$0{,}8\cos\!\left(\frac{\pi t}{6}\right)+1 = 1{,}4 \;\Rightarrow\; \cos\!\left(\frac{\pi t}{6}\right) = 0{,}5$$

As soluções gerais de $\cos\theta = 0{,}5$ são $\theta = \dfrac{\pi}{3}+2k\pi$ ou $\theta = -\dfrac{\pi}{3}+2k\pi$ (equivalente a $\dfrac{5\pi}{3}+2k\pi$).

Fazendo $\theta = \dfrac{\pi t}{6}$, temos $t = \dfrac{6\theta}{\pi}$:

\begin{itemize}[nosep,leftmargin=*]
\item Para $\theta=\dfrac{\pi}{3}$: $t = 6\cdot\dfrac{1}{3} = 2$
\item Para $\theta=\dfrac{5\pi}{3}$: $t = 6\cdot\dfrac{5}{3} = 10$
\end{itemize}

Dentro de $0\le t<12$, as soluções são $t=2$ e $t=10$.

\textbf{Análise de crescimento/decrescimento}

Como $h(0)=1{,}8$ é máximo e $h(6)=0{,}2$ é mínimo, no intervalo $(0,6)$ a função cosseno está descendo de seu valor máximo ao mínimo — ou seja, $h(t)$ é \textbf{decrescente} nesse trecho. Logo, em $t=2$ (que está entre 0 e 6), a maré está \textbf{descendo}.

No intervalo $(6,12)$, a função retorna do mínimo ($h(6)=0{,}2$) ao próximo máximo ($h(12)=1{,}8$), portanto $h(t)$ é \textbf{crescente} nesse trecho. Logo, em $t=10$ (entre 6 e 12), a maré está \textbf{subindo}.

\textbf{Conclusão:} às 02h00 a maré está a 1,4 m e descendo; às 10h00 a maré está a 1,4 m e subindo.

\bigskip\noindent\textbf{Q087 · múltipla escolha · média · aplicar · teórica · gabarito conferido}\par\nopagebreak
Considere a progressão aritmética $(a_n)$ com primeiro termo $a_1 = 7$ e razão $r = 4$: $(7, 11, 15, 19, 23, \dots)$. Essa PA pode ser vista como a restrição da função afim $f(x) = 4x + 3$ ao conjunto dos números naturais não nulos, de modo que $a_n = f(n)$ para todo $n \geq 1$ (o gráfico da PA é o conjunto de pontos discretos $(n, a_n)$ sobre a reta que representa $f$).

Qual é o menor termo dessa PA que é maior do que 100?
\begin{enumerate}[label=(\alph*)]
\item 100
\item 99
\item 103
\item 107
\end{enumerate}

\textbf{Gabarito.} 103

\textbf{Resolução.} \textbf{Passo 1 — Associar a PA à função afim.}
O termo geral da PA é $a_n = a_1 + (n-1)r = 7 + 4(n-1) = 4n + 3$. Logo, a PA corresponde à função afim $f(x) = 4x+3$ restrita aos naturais não nulos, isto é, $a_n = f(n)$, com $n \in \{1, 2, 3, \dots\}$.

\textbf{Passo 2 — Resolver a condição sobre os números reais.}
Queremos o menor termo maior que 100, ou seja, o menor $n$ tal que $f(n) > 100$. Primeiro resolvemos a inequação como se $x$ fosse real:
$$4x + 3 > 100 \implies 4x > 97 \implies x > 24{,}25.$$

\textbf{Passo 3 — Voltar ao domínio discreto da PA.}
Como $n$ representa a posição de um termo, ele precisa ser um número \textbf{natural}. O menor número natural maior que $24{,}25$ é $n = 25$ (e não $24$, que não satisfaz a inequação, nem $24{,}25$, que não é índice de termo algum).

\textbf{Passo 4 — Calcular o termo.}
$$a_{25} = f(25) = 4(25) + 3 = 103.$$

Como verificação: $a_{24} = 4(24)+3 = 99 < 100$ e $a_{25} = 103 > 100$, confirmando que $a_{25}=103$ é o menor termo da PA maior que 100.

\bigskip\noindent\textbf{Q061 · discursiva · fácil · aplicar · aplicada · gabarito conferido · corrigida após a avaliação docente}\par\nopagebreak
Um biólogo inicia uma cultura de bactérias com 200 indivíduos. A cada hora que passa, a quantidade de bactérias na cultura dobra em relação à hora anterior. Seja $a_n$ o número de bactérias existentes ao final da $n$-ésima hora de observação, com $n = 1, 2, 3, \dots$ (isto é, $a_1$ é a quantidade ao final da 1ª hora, $a_2$ ao final da 2ª hora, e assim por diante).

a) Mostre que a sequência $(a_n)$ é uma progressão geométrica, indicando o primeiro termo e a razão.

b) Escreva a lei $f(n)$ de uma função exponencial que descreva exatamente essa progressão geométrica, especificando qual é o domínio dessa função e por que ele deve ser desse tipo (e não todos os números reais).

c) Usando a fórmula do termo geral da PG (ou, equivalentemente, a função $f(n)$ obtida no item b), determine quantas bactérias haverá ao final da 6ª hora de observação.

\textbf{Gabarito.} a) PG de razão $q=2$ e $a_1=400$ (as 200 iniciais já dobraram ao final da 1ª hora); b) $f(n)=400\cdot 2^{n-1}=200\cdot 2^{n}$, com domínio $D=\{1,2,3,\dots\}$ (naturais, pois $n$ conta horas); c) $a_6 = f(6) = 12800$ bactérias.

\textbf{Resolução.} \textbf{a) Identificando a PG}

Como a quantidade dobra a cada hora, temos $a_{n+1} = 2\cdot a_n$ para todo $n \geq 1$. Isso significa que a razão entre termos consecutivos é constante e igual a $q = 2$. Logo, $(a_n)$ é uma progressão geométrica de razão $q=2$.

O primeiro termo exige atenção: as 200 bactérias são a quantidade \textbf{no início} da observação, e $a_1$ é a quantidade \textbf{ao final} da 1ª hora, quando ela já dobrou uma vez. Assim, $a_1 = 2 \cdot 200 = 400$.

\textbf{b) Associando a PG a uma função exponencial}

O termo geral da PG é dado por
$$a_n = a_1 \cdot q^{\,n-1} = 400 \cdot 2^{\,n-1} = 200 \cdot 2^{\,n}.$$

Essa expressão é exatamente a lei de uma função exponencial
$$f(n) = 200 \cdot 2^{\,n},$$
em que a base da exponencial ($2$) coincide com a razão da PG e o coeficiente ($200$) é a quantidade inicial — o valor que a curva assumiria em $n = 0$, instante que não pertence à sequência.

Como $n$ representa o número da hora de observação (1ª hora, 2ª hora, etc.), $n$ só pode assumir valores inteiros positivos: $n \in \{1, 2, 3, \dots\}$. Assim, o domínio de $f$ é o conjunto dos números naturais (a partir de 1), $D = \mathbb{N}^* = \{1,2,3,\dots\}$, e \textbf{não} todo o conjunto dos reais — não faz sentido, por exemplo, falar em "1,5 hora de observação" nesse contexto discreto de contagem de termos da sequência. Por isso $f$ é uma função exponencial de domínio discreto que apenas "amostra" a curva exponencial contínua $y = 200\cdot 2^{x}$ nos pontos inteiros.

\textbf{c) Cálculo para a 6ª hora}

Usando $n=6$:
$$f(6) = a_6 = 200 \cdot 2^{6} = 200 \cdot 64 = 12800.$$

Logo, ao final da 6ª hora haverá \textbf{12800 bactérias}.

\endgroup


\chapter{Instrumento e respostas da primeira rodada de avaliação}
\label{ap:instrumento}

Cada professor recebeu um caderno impresso com 21 ou 22 questões. Cada item trazia o enunciado, as
alternativas quando houvesse, o gabarito proposto e a habilidade da BNCC declarada, e era seguido das
quatro perguntas abaixo.

\begin{enumerate}
\item \textbf{Esta questão tem erro matemático?} Não. Sim, e o erro está em \underline{\hspace{4cm}}.
  Não sei dizer.
\item \textbf{Nível de dificuldade para os seus alunos.} Fácil. Média. Difícil. Fora do alcance da
  turma.
\item \textbf{Decisão de uso.} Aceita. Aceita com ajuste. Recusada.
\item \textbf{Por quê?} Resposta obrigatória quando a decisão não for \emph{aceita}.
\end{enumerate}

A instrução de abertura informava que a amostra tinha qualidade variável e podia conter erros
matemáticos, e que não se avaliava o participante e sim o material. Ao final do caderno, três
perguntas abertas completavam o instrumento.

\begin{enumerate}
\item Você usaria um sistema como este no seu planejamento? Por quê?
\item O que faltou nas questões ou no material?
\item O que atrapalhou?
\end{enumerate}

A Tabela~\ref{tab:apendice-respostas} reúne as decisões de uso de todos os avaliadores sobre todas
as questões. As colunas indicam o código da habilidade sem o prefixo \texttt{EM13MAT}, o formato,
discursiva ou de múltipla escolha, a dificuldade declarada, que pode ser fácil, média ou difícil, e
a garantia obtida, em que C indica gabarito conferido, P conferido em parte e S sem conferência
automática. Nas colunas dos avaliadores, A indica aceita, J aceita com ajuste, R recusada e o
traço uma decisão em branco, enquanto a célula vazia indica que a questão não estava no lote daquele
avaliador. A marca $e$ assinala as respostas em que o avaliador indicou erro matemático, nenhum deles
confirmado, e o asterisco assinala a Q061, a única questão com erro, que nenhum avaliador apontou.

% Gerado a partir de sistema/medicoes/acervo-consolidado/avaliacao/ (_chave.csv e
% respostas/transcricao-avaliador-*.csv). Não editar à mão: regerar se os dados mudarem.
{\small
\begin{longtable}{llllcccccc}
\caption{Decisões de uso na primeira rodada, por questão}\label{tab:apendice-respostas}\\
\toprule
Item & Hab. & Formato & Decl. & Gar. & Av.~1 & Av.~2 & Av.~3 & Av.~4 & Av.~5 \\
\midrule
\endfirsthead
\multicolumn{10}{l}{\footnotesize\emph{continuação}}\\
\toprule
Item & Hab. & Formato & Decl. & Gar. & Av.~1 & Av.~2 & Av.~3 & Av.~4 & Av.~5 \\
\midrule
\endhead
\bottomrule
\endfoot
Q001 & 302 & Disc. & F & C &  & A &  &  &  \\
Q002 & 405 & Disc. & F & P &  &  &  &  & -- \\
Q003 & 302 & Disc. & M & C & A &  &  &  &  \\
Q004 & 405 & Disc. & M & P & A &  &  &  &  \\
Q005 & 302 & Disc. & D & P &  & J &  &  &  \\
Q006 & 405 & Disc. & D & P &  &  &  & A &  \\
Q007 & 302 & ME & F & C & A & A & A & A & J \\
Q008 & 405 & ME & F & P &  & J &  &  &  \\
Q009 & 302 & ME & M & C &  & A &  &  &  \\
Q010 & 405 & ME & M & P &  &  & R &  &  \\
Q011 & 302 & ME & D & C &  &  & J &  &  \\
Q012 & 405 & ME & D & S & A &  &  &  &  \\
Q013 & 403 & Disc. & F & P &  & J &  &  &  \\
Q014 & 501 & Disc. & F & C &  &  &  &  & A \\
Q015 & 403 & Disc. & M & C & A & R & R & A & J \\
Q016 & 501 & Disc. & M & C &  & A &  &  &  \\
Q017 & 403 & Disc. & D & P &  &  &  &  & J \\
Q018 & 501 & Disc. & D & C &  &  &  & A &  \\
Q019 & 403 & ME & F & P &  &  &  & A &  \\
Q020 & 501 & ME & F & C & A &  &  &  &  \\
Q021 & 403 & ME & M & P &  & J &  &  &  \\
Q022 & 501 & ME & M & C &  &  &  &  & A \\
Q023 & 403 & ME & D & C &  &  &  & A$^{e}$ &  \\
Q024 & 501 & ME & D & C & A & A & A & A & A \\
Q025 & 401 & Disc. & F & C & A &  &  &  &  \\
Q026 & 404 & Disc. & F & C &  &  &  &  & R \\
Q027 & 502 & Disc. & F & C &  &  & A &  &  \\
Q028 & 401 & Disc. & M & C &  &  & A &  &  \\
Q029 & 404 & Disc. & M & C &  &  & R &  &  \\
Q030 & 502 & Disc. & M & C &  &  &  & A &  \\
Q031 & 401 & Disc. & D & C &  &  & J &  &  \\
Q032 & 404 & Disc. & D & C &  & R &  &  &  \\
Q033 & 502 & Disc. & D & C &  &  & A &  &  \\
Q034 & 401 & ME & F & C & A & A & J & A & J \\
Q035 & 404 & ME & F & C &  & J &  &  &  \\
Q036 & 502 & ME & F & C &  &  &  & A &  \\
Q037 & 401 & ME & M & C & A &  &  &  &  \\
Q038 & 404 & ME & M & C &  &  & R &  &  \\
Q039 & 502 & ME & M & C & A &  &  &  &  \\
Q040 & 401 & ME & D & C &  &  &  &  & A \\
Q041 & 404 & ME & D & C &  &  &  & R &  \\
Q042 & 502 & ME & D & C & J &  &  &  &  \\
Q044 & 306 & Disc. & F & C &  &  & J &  &  \\
Q045 & 507 & Disc. & F & C & A &  &  &  &  \\
Q046 & 304 & Disc. & M & C &  &  & A &  &  \\
Q047 & 402 & Disc. & M & C &  &  &  &  & J \\
Q048 & 508 & Disc. & M & C &  & R &  &  &  \\
Q049 & 305 & Disc. & D & C &  &  &  &  & R \\
Q050 & 503 & Disc. & D & P &  &  &  & A &  \\
Q052 & 306 & ME & F & C & -- &  &  &  &  \\
Q053 & 507 & ME & F & C & -- &  &  &  &  \\
Q054 & 304 & ME & M & C &  &  &  &  & A \\
Q055 & 402 & ME & M & C &  & J &  &  &  \\
Q056 & 508 & ME & M & C &  &  &  &  & A \\
Q057 & 305 & ME & D & C &  & R &  &  &  \\
Q058 & 503 & ME & D & P &  &  &  & A &  \\
Q061$^{*}$ & 508 & Disc. & F & C & A &  &  &  &  \\
Q062 & 305 & Disc. & M & C &  &  & -- &  &  \\
Q063 & 503 & Disc. & M & C &  &  & A &  &  \\
Q064 & 303 & Disc. & D & C &  & R &  &  &  \\
Q065 & 306 & Disc. & D & C & A & R & R & J & R \\
Q066 & 507 & Disc. & D & C &  &  &  & A &  \\
Q067 & 304 & ME & F & C &  &  & J &  &  \\
Q068 & 402 & ME & F & C &  &  &  & J &  \\
Q069 & 508 & ME & F & C &  &  & R &  &  \\
Q070 & 305 & ME & M & C &  &  &  & J &  \\
Q071 & 503 & ME & M & C & -- &  &  &  &  \\
Q072 & 303 & ME & D & C &  & J &  &  &  \\
Q073 & 306 & ME & D & C &  &  &  &  & R$^{e}$ \\
Q074 & 507 & ME & D & C &  & A$^{e}$ &  &  &  \\
Q075 & 305 & Disc. & F & C &  &  &  &  & R \\
Q076 & 503 & Disc. & F & C &  &  &  & A &  \\
Q077 & 303 & Disc. & M & C & A &  &  &  &  \\
Q078 & 306 & Disc. & M & C &  &  &  &  & R \\
Q079 & 507 & Disc. & M & C &  &  &  & A &  \\
Q080 & 304 & Disc. & D & C &  &  &  & A &  \\
Q081 & 402 & Disc. & D & C &  &  & J &  &  \\
Q082 & 508 & Disc. & D & C &  &  & J &  &  \\
Q084 & 503 & ME & F & C &  &  &  &  & A$^{e}$ \\
Q085 & 303 & ME & M & C & A &  &  &  &  \\
Q086 & 306 & ME & M & C & A &  &  &  &  \\
Q087 & 507 & ME & M & C & A & A & A & A & A \\
Q088 & 304 & ME & D & C &  &  &  &  & A \\
Q089 & 402 & ME & D & C &  &  &  & A &  \\
Q090 & 508 & ME & D & C &  & A &  &  &  \\
\end{longtable}
}


\chapter{Termo de Consentimento Livre e Esclarecido}
\label{ap:tcle}

% CONFERIR: este é o modelo do termo, redigido a partir do que o rascunho
% da aluna descreve (objetivos, natureza do material, ausência de respostas
% certas ou erradas, anonimato). Substituir pelo texto efetivamente assinado
% pelos participantes, se ele for diferente.

\begin{center}
\textbf{TERMO DE CONSENTIMENTO LIVRE E ESCLARECIDO}
\end{center}

Você está sendo convidado a participar, como avaliador, da pesquisa \emph{Sistema multiagente para
geração assistida de questões de Matemática no Ensino Médio}, desenvolvida por Isabella Dias
Ribeiro dos Santos no Mestrado Profissional em Matemática em Rede Nacional (PROFMAT) da Universidade
Federal de Mato Grosso, sob orientação do Prof. Dr. Moiseis dos Santos Cecconello.

\textbf{Objetivo.} A pesquisa investiga um sistema computacional que gera questões de Matemática
do Ensino Médio com o auxílio de inteligência artificial, confere automaticamente os cálculos e
avalia a qualidade didática das questões. A sua participação consiste em examinar um conjunto de
questões e registrar, para cada uma, se ela contém erro matemático, qual o seu nível de dificuldade
para os seus alunos e se você a usaria em sala de aula, respondendo ao final a três perguntas
abertas.

\textbf{Duração e riscos.} A atividade leva cerca de 75 minutos e pode ser feita no momento que
lhe for mais conveniente. Os riscos são mínimos e se limitam ao tempo dedicado e a um eventual
desconforto ao emitir julgamentos sobre o material. Não há respostas certas ou erradas, pois o
objeto de avaliação são as questões e não o participante.

\textbf{Benefícios.} A participação contribui para o aprimoramento de um produto educacional de
acesso livre destinado a professores de Matemática. Ao final, você pode solicitar à pesquisadora a
indicação das questões que continham erro e a explicação de cada um.

\textbf{Sigilo e voluntariedade.} A sua identidade não será divulgada em nenhuma hipótese. Os
avaliadores são identificados apenas por numeração, sem relação com nenhuma característica pessoal,
e nenhum dado individual identificável será publicado. A participação é voluntária e não
remunerada, e você pode desistir a qualquer momento, sem qualquer prejuízo, bastando comunicar a
pesquisadora.

Declaro que fui informado dos objetivos e procedimentos da pesquisa, que tive a oportunidade de
esclarecer as minhas dúvidas e que concordo em participar voluntariamente.

\vspace{1.5cm}
\noindent\rule{8cm}{0.4pt}\\
Assinatura do participante \hfill Data: \underline{\hspace{3cm}}

\vspace{1cm}
\noindent\rule{8cm}{0.4pt}\\
Isabella Dias Ribeiro dos Santos, pesquisadora responsável

\chapter{Repositório e código-fonte}
\label{ap:repositorio}

O código do sistema, os \emph{prompts}, o catálogo de habilidades, os registros das execuções com
provedor real, o banco curado e a fonte desta dissertação estão reunidos em um único repositório
público, disponível em \url{https://github.com/1moi6/isabela}. A organização da pasta
\texttt{sistema/} foi apresentada na Seção~\ref{sec:repositorio}, e as demais pastas cumprem os
papéis descritos a seguir.

\begin{itemize}
\item \texttt{banco\_curado/} contém o banco de 85 questões em JSON e em PDF, gerado a partir do
  acervo por \texttt{sistema/montar\_banco.py}.
\item \texttt{sistema/medicoes/} guarda cada execução com provedor real em uma pasta própria, com o
  registro completo dos ciclos, uma descrição da execução com data, modelo e versão do código e um
  texto que resume o que ela mostrou. O acervo usado nesta dissertação está em
  \texttt{acervo-consolidado/}, que inclui também o material entregue aos avaliadores e as
  transcrições das respostas.
\item \texttt{docs/} contém a página do professor, que é servida tanto pelo programa local quanto
  pela publicação do repositório na web.
\item \texttt{dissertacao/} contém a fonte \LaTeX{} deste texto.
\end{itemize}

Para obter o sistema basta copiar o repositório. O professor que não programa encontra em
\texttt{sistema/LEIA-ME.txt} o roteiro de instalação por duplo clique descrito na
Seção~\ref{sec:produto-sistema}. Quem desejar conferir a implementação pode instalar o pacote com
as dependências de desenvolvimento e executar a suíte de testes, que não exige chave de API.

\begin{verbatim}
cd sistema
pip install -e ".[dev]"
python -m pytest
\end{verbatim}

Os planos que orientaram as etapas ainda não executadas, a saber o painel formal e a comparação
entre modelos, estão na raiz do repositório, em \texttt{plano\_experimento\_painel.md} e
\texttt{plano\_teste\_modelos.md}.

\chapter{Manual do professor: uso e limitações}
\label{ap:manual}

Este apêndice reproduz a segunda parte do manual do professor, que complementa o guia de instalação
distribuído com o programa. Ela foi escrita para ser lida por quem nunca usou uma ferramenta desse
tipo.

\section*{Como pedir uma questão}

O formulário começa pela \textbf{habilidade da BNCC}, e não pelo assunto, porque é a habilidade que
diz o que a questão precisa exigir do aluno. Escolhida a habilidade, o sistema mostra apenas os temas
que ela cobre. Em algumas habilidades a combinação de temas é obrigatória, como na
\texttt{EM13MAT507}, que associa progressão aritmética e função afim, e nesse caso o sistema não
aceita pedir só um deles.

O \textbf{nível cognitivo} diz que tipo de pensamento a questão deve pedir, de lembrar e entender até
analisar, avaliar e criar. O sistema sugere os níveis que combinam com o verbo da habilidade, mas
você pode escolher outro. A \textbf{dificuldade} tem quatro degraus. Use o nível \emph{básico} para
turmas que ainda não dominam os pré-requisitos da habilidade, pois ele pede um único passo, números
pequenos e enunciado curto. Evite combinar o nível básico com o processo \emph{criar}, porque uma
questão de um só passo raramente exige criação, e o sistema avisa quando isso acontece. A
\textbf{natureza} define se a questão é teórica ou aplicada a uma situação do cotidiano, e o
\textbf{formato} se ela é discursiva ou de múltipla escolha. O \textbf{contexto} e as
\textbf{restrições} são opcionais e servem para aproximar a questão da sua turma, por exemplo pedindo
um contexto de esportes ou apenas números inteiros.

\section*{Três fluxos de trabalho recomendados}

O primeiro fluxo serve para \textbf{montar uma avaliação}. Gere de três a cinco questões por
habilidade que pretende avaliar, leia a trilha de cada uma, guarde no banco as que lhe parecerem
adequadas e registre a sua decisão. Depois selecione as questões na área do banco e exporte a lista
nas duas versões, a do aluno para imprimir e a do professor para a correção.

O segundo fluxo serve para \textbf{construir um acervo pessoal ao longo do bimestre}. A cada
conteúdo trabalhado, gere algumas questões em níveis diferentes e registre a sua avaliação. Com o
tempo o seu banco passa a refletir a sua turma, e os comentários que você deixou mostram o que
precisou ser ajustado e por quê.

O terceiro fluxo serve para \textbf{diferenciar uma lista por nível}. Peça a mesma habilidade em
dificuldades diferentes, do básico ao difícil, e monte listas distintas para grupos de alunos em
momentos diferentes da aprendizagem. Quem não instalou o sistema pode fazer o mesmo a partir do
banco curado em PDF, que está organizado por habilidade.

\section*{Limitações conhecidas e como contorná-las}

\textbf{O sistema erra, e o seu julgamento continua sendo indispensável.} Toda questão, inclusive
as aprovadas, deve ser lida e resolvida por você antes de chegar à sala de aula. O sistema reduz o
trabalho de produzir a questão e mostra o que foi conferido, mas não substitui a sua revisão.

\textbf{O que significa cada garantia.} \emph{Gabarito conferido} quer dizer que o computador
refez as contas de forma independente e chegou ao mesmo resultado do gabarito. \emph{Conferido em
parte} quer dizer que só uma parte do que a questão afirma pôde ser refeita, e o restante depende da
sua leitura. \emph{Sem conferência automática} quer dizer que nada foi refeito, e por isso essas
questões merecem a sua conferência antes das demais.

\textbf{Mesmo o gabarito conferido pode estar errado.} O computador confere as contas a partir da
forma como o próprio sistema entendeu o enunciado. Se o sistema entendeu mal o enunciado, a conta
confere e o gabarito continua errado. Isso aconteceu uma vez no banco, numa questão em que a
quantidade inicial de bactérias foi tomada como o primeiro termo da progressão quando o enunciado
dizia outra coisa. Ao revisar, confira sobretudo se o gabarito responde ao que o enunciado de fato
pergunta.

\textbf{A dificuldade é a de um aluno que cumpre a habilidade no ano previsto.} Na avaliação com
professores da rede pública, a maior parte das questões foi julgada mais difícil do que o declarado.
Se a sua turma está aquém da habilidade, comece pelo nível básico e ajuste a partir daí.

\textbf{A qualidade depende do provedor escolhido.} As questões do banco foram geradas com um modelo
comercial de grande porte. Modelos menores e modelos abertos executados no próprio computador custam
menos ou nada, mas a qualidade das questões produzidas por eles não foi medida nesta pesquisa, e a
revisão deve ser ainda mais cuidadosa.

\textbf{O sistema produz apenas texto.} Ainda não há figuras nem gráficos. Nas habilidades que pedem
passar da forma algébrica para o gráfico, acrescente você mesmo o esboço quando ele for necessário.

\textbf{Contextos podem se repetir.} Ao gerar muitas questões da mesma habilidade, o sistema tende a
reaproveitar situações, como culturas de bactérias ou corridas por aplicativo. Use o campo de
contexto para variar as situações.

\textbf{A redação da habilidade segue o texto oficial da BNCC.} Se a sua rede usa um documento
curricular próprio, a redação pode ser diferente da que aparece no sistema. Na dúvida, confira o
código no documento da Base.
